\section{Methodology}
\label{sec:method}

ReAD treats Problem~\ref{prob:budget_alloc} as state-dependent budget allocation
over capabilities. At each step, it estimates task requirements, generates
capability-targeted teacher supervision, and uses an uncertainty-aware contextual
bandit to choose the next token allocation. As the student profile changes, ReAD
can shift budget away from saturated or high-spillover capabilities.
\rev{ReAD is \emph{off-policy} in the sense of \citet{agarwal2024gkd}: the
student trains on teacher completions to generated prompts and never samples
its own training sequences. The allocator is orthogonal to the inner objective,
and on-policy losses~\citep{gu2024minillm,ko2024distillm} could replace the
token-level loss in Section~\ref{sec:method:bandit} unchanged.}

\subsection{Identifying Task-Essential Capabilities}
\label{sec:method:essential}

\paragraph{Task requirement vector.}
For each downstream task $\tau$, ReAD constructs a task card
$\mathcal{D}^{\mathrm{spec}}_\tau$ from non-test data, containing a short task
description, input/output format, evaluation target, and a few representative
exemplars, with the test split reserved only for final reporting. From this
task card, ReAD estimates a requirement vector
$\mathbf{r}_\tau\in\Delta_{|\mathcal{C}|}$, where $r_{\tau,c}$ indicates how
useful capability $c$ is for improving the task utility $U_\tau(\cdot)$. We use
$\mathbf{r}_\tau$ as an actionable allocation prior: high-mass capabilities are
prioritized for improvement, and drops on these capabilities are treated as
harmful spillover. Since token allocations cannot be negative,
$\mathbf{r}_\tau$ is constrained to the simplex; negative sensitivities are
captured instead by signed capability changes and the spillover penalty in
Section~\ref{sec:method:bandit}.

\paragraph{Requirement identifier.}
We learn a lightweight identifier $g_\phi$ that maps the task card to the requirement vector $\mathbf{r}_\tau = g_\phi(\mathcal{D}^{\mathrm{spec}}_\tau) \in \Delta_{|\mathcal{C}|}$.
% \begin{equation}
% \mathbf{r}_\tau = g_\phi(\mathcal{D}^{\mathrm{spec}}_\tau) \in \Delta_{|\mathcal{C}|}.
% \label{eq:req_id}
% \end{equation}
The identifier is a small Transformer encoder over the task card, followed by a two-layer MLP and a softmax head. Since $\mathbf{r}_\tau$ is not directly observed, we supervise it through low-budget interventions in capability space.

\paragraph{Local supervision signal.}
Let $\mathbf{s}(S)\in\mathbb{R}^{|\mathcal{C}|}$ denote the measured capability profile of a student $S$, and write $U_\tau(S)=F_\tau(\mathbf{s}(S))$. For a small intervention that changes the profile from $\mathbf{s}$ to $\mathbf{s}+\Delta\mathbf{s}$,
\begin{equation}
F_\tau(\mathbf{s}+\Delta\mathbf{s})
=
F_\tau(\mathbf{s})+
\nabla_{\mathbf{s}}F_\tau(\mathbf{s})^\top\Delta\mathbf{s}
+o(\|\Delta\mathbf{s}\|),
\label{eq:taylor}
\end{equation}
so the utility change is locally approximated by
\begin{equation}
\Delta U_\tau
=
U_\tau(S')-U_\tau(S)
\approx
\nabla_{\mathbf{s}}F_\tau(\mathbf{s})^\top\Delta\mathbf{s}.
\label{eq:local_linear}
\end{equation}
Thus, $\mathbf{r}_\tau$ is trained as a simplex-constrained surrogate for the beneficial part of the local utility sensitivity, e.g., $\mathbf{r}_\tau\approx\Pi_\Delta([\nabla_{\mathbf{s}}F_\tau(\mathbf{s})]_+)$.

\paragraph{Offline pretraining.}
We construct an intervention set
$\mathcal{Z}=\{(\mathcal{D}^{\mathrm{spec}}_\tau,\Delta\mathbf{s}^{(m)}_\tau,\Delta U^{(m)}_\tau)\}_{\tau,m}$ from auxiliary tasks and non-test splits. For each task $\tau$, we sample sparse allocation vectors $\mathbf{w}^{(m)}\in\Delta_{|\mathcal{C}|}$ by choosing a support set $\mathcal{I}^{(m)}\subseteq\mathcal{C}$, sampling nonzero weights from a Dirichlet distribution, and renormalizing them. A short probe distillation under a small budget $b_{\mathrm{probe}}$ produces an intervened student $S^{(m)}$, from which we measure
\begin{equation}
\Delta U^{(m)}_\tau := U_\tau(S^{(m)})-U_\tau(S_0),
\qquad
\Delta\mathbf{s}^{(m)}_\tau := \mathbf{s}(S^{(m)})-\mathbf{s}(S_0).
\label{eq:probe_deltas}
\end{equation}
We pretrain $g_\phi$ by predicting the observed utility change from the induced capability change:
\begin{equation}
\min_\phi
\sum_{(\tau,m)}
\Big(\Delta U^{(m)}_\tau
- g_\phi(\mathcal{D}^{\mathrm{spec}}_\tau)^\top\Delta\mathbf{s}^{(m)}_\tau\Big)^2
+\lambda_{\mathrm{ent}}\,\mathcal{H}\!\left(g_\phi(\mathcal{D}^{\mathrm{spec}}_\tau)\right).
\label{eq:req_pretrain}
\end{equation}
The positive entropy penalty discourages near-uniform requirement vectors and makes the allocation prior more discriminative. At deployment time, ReAD computes $\mathbf{r}_\tau$ once and reuses it throughout the distillation run. Since the identifier is trained once and reused across tasks, it is not counted as part of the per-run student distillation budget.

\subsection{On-the-Fly Capability-Targeted Data Generation}
\label{sec:method-generator}

\paragraph{Capability-conditioned templates.}
Given $\mathbf{r}_\tau$, ReAD generates training examples on the fly from a capability-labeled template library. For each capability $c\in\mathcal{C}$, a template set $\mathcal{P}_c$ specifies an instruction scaffold, typed slots, and output-format constraints. Slot values are filled using deterministic seeds. The generator controls the form of teacher supervision

\paragraph{Difficulty scoring.}
For prompt $x$ in capability $c$, we compute a deterministic difficulty score
\begin{equation}
d_c(x)=\frac{1}{|\mathcal{F}_c|}\sum_{f\in\mathcal{F}_c}
\frac{f(x)-\min_{x'\in\mathcal{Q}_c} f(x')}
{\max_{x'\in\mathcal{Q}_c} f(x')-\min_{x'\in\mathcal{Q}_c} f(x')+\epsilon},
\label{eq:difficulty_score}
\end{equation}
where $\mathcal{F}_c$ is the set of capability-specific control factors and $\mathcal{Q}_c$ is a calibration prompt pool for capability $c$. For example, reasoning templates use factors such as the number of constraints, reasoning hops, and symbolic complexity; code templates use the number of required functions and tests; steerability templates use the number of rules and schema fields. We split each capability's calibration pool into easy, medium, and hard buckets by tertiles of $d_c(x)$.

\paragraph{Curriculum and sampling.}
At decision step $t$, ReAD samples a capability label according to the selected allocation $\mathbf{w}_t$, instantiates a prompt from the corresponding template set, queries the teacher $T$ for a completion $y$, and immediately distills on $(x,y)$. The curriculum only adjusts the within-capability difficulty mix: early steps emphasize easy and medium examples, while later steps increase the probability of hard examples. It does not choose the capability allocation. We also track recent template frequencies and cap each template's share so that no single scaffold dominates.

\subsection{Contextual Bandit for Capability Allocation}
\label{sec:method:bandit}

\paragraph{Decision loop.}
ReAD runs for $T_{\mathrm{step}}=20$ decision steps\rev{, with the bandit update, data resampling, and proxy refresh at every step (Section~\ref{subsec:exp_setup})}. At step $t$, the context summarizes the task requirement, current student profile, remaining budget, and recent allocation history:
\begin{equation}
\mathbf{x}_t=[\mathbf{r}_\tau;\mathbf{s}^{\mathrm{probe}}(S_t);b_t;\boldsymbol{\rho}_t],
\label{eq:bandit_context}
\end{equation}
where $\boldsymbol{\rho}_t$ stores a short record of recent allocations and observed gains. The profile $\mathbf{s}^{\mathrm{probe}}(S_t)$ is computed with a small fixed monitoring suite, while full held-out benchmark evaluation is used only for final reporting and coarse calibration.

\paragraph{Candidate allocations and distillation.}
At each step, ReAD selects an allocation $\mathbf{w}_t$ from a finite candidate set $\mathcal{A}(\tau)$ consisting of local perturbations around the previous allocation and sparse actions concentrated on the top-$k$ capabilities under $\mathbf{r}_\tau$. Weights are chosen from a fixed grid and renormalized to sum to one. Given $\mathbf{w}_t$, ReAD samples capability-targeted data and updates the student by minimizing $\mathcal{L}_{\mathrm{distill}}(\theta_t)
=-\mathbb{E}_{(x,y)}\sum_{j=1}^{|y|}
\log p_{\theta_t}(y_j\mid x,y_{<j})$ to produce $S_{t+1}$.
% \begin{equation}
% \mathcal{L}_{\mathrm{distill}}(\theta_t)
% =-\mathbb{E}_{(x,y)}\sum_{j=1}^{|y|}
% \log p_{\theta_t}(y_j\mid x,y_{<j}),
% \label{eq:distill_loss}
% \end{equation}

\paragraph{Proxy reward.}
After the update, ReAD measures the probe-profile change
$\Delta\mathbf{s}^{\mathrm{probe}}_t=\mathbf{s}^{\mathrm{probe}}(S_{t+1})-\mathbf{s}^{\mathrm{probe}}(S_t)$ and forms
\begin{equation}
\widehat{R}_t
=\mathbf{r}_\tau^\top\Delta\mathbf{s}^{\mathrm{probe}}_t
-\beta\,\mathrm{Spill}_t
-\lambda\,\mathrm{cost}_t, \quad\ \mathrm{Spill}_t=
\sum_{c\in\mathcal{C}_{\mathrm{ess}}}
 r_{\tau,c}\,[-\Delta s^{\mathrm{probe}}_{t,c}]_+
\label{eq:reward}
\end{equation}
Here $\mathcal{C}_{\mathrm{ess}}$ contains the top-$k$ capabilities under $\mathbf{r}_\tau$, and $\mathrm{cost}_t$ is the token budget consumed at the step. The first term rewards task-aligned capability gains, the second penalizes regressions on task-essential capabilities, and the third enforces budget awareness.

\paragraph{Reward model and UCB allocation.}
We append $(\mathbf{x}_t,\mathbf{w}_t,\widehat{R}_t)$ to a history buffer and train an ensemble of $J$ two-layer MLP reward regressors $\{h_{\eta_j}\}_{j=1}^J$. Each regressor maps $(\mathbf{x},\mathbf{w})$ to a scalar next-step reward and is trained on a bootstrap resample of the 10\% profiling split plus logged transitions from earlier steps, yielding about $3.2$k--$4.9$k examples per task in our experiments. For a candidate allocation $\mathbf{w}$, the ensemble mean and uncertainty are
\begin{equation}
\mu(\mathbf{x}_t,\mathbf{w})=\frac{1}{J}\sum_{j=1}^J h_{\eta_j}(\mathbf{x}_t,\mathbf{w}),
\quad
\sigma(\mathbf{x}_t,\mathbf{w})=
\sqrt{\frac{1}{J-1}\sum_{j=1}^J\big(h_{\eta_j}(\mathbf{x}_t,\mathbf{w})-\mu(\mathbf{x}_t,\mathbf{w})\big)^2}.
\label{eq:ensemble_mu_sigma}
\end{equation}
ReAD then selects the next allocation using an upper-confidence rule,
\begin{equation}
\mathbf{w}_{t+1}
=\arg\max_{\mathbf{w}\in\mathcal{A}(\tau)}
\mu(\mathbf{x}_{t+1},\mathbf{w})+\kappa\,\sigma(\mathbf{x}_{t+1},\mathbf{w}),
\label{eq:ucb}
\end{equation}
with $b_{t+1}=b_t-\mathrm{cost}_t$. At coarse development checkpoints, we fit an affine calibration from cumulative proxy reward to observed utility change and reduce exploration if the development utility stagnates. Unlike static mixture-regression approaches such as RegMix~\citep{liu2024regmix}, ReAD repeatedly re-estimates the best next allocation from the current student state and remaining budget.
